\subsection{Where Reported Wins over SMOTE Come From}
\label{sec:data-wins}
% Rewritten 29 Oct 2026 (Claude; Toon delegated every remaining decision on 29 Oct).
% Win/loss definition decided by Toon (28 Oct 2026): every metric the paper reports
% with a SMOTE column, classifier and condition fixed by rule A, printed values
% compared in the metric's direction, equal printed values = tie. Numbers:
% out/winloss_allmetric_summary.json, out/hypotheses.json, out/difficulty_outcomes.json.
At the paper level, 44 of the 73 papers use plain SMOTE as a baseline,
13 use SMOTE inside a multiclass or ensemble scheme, and 16 do not compare with
SMOTE at all. Fewer papers give scores a reader can pair: 29 papers print
per-dataset scores for the proposed method and plain SMOTE in the same table.
For these papers we compared the printed values on every metric that the table
reports for both methods, with the classifier and condition fixed as in
Section~\ref{sec:data-method}; a win means the proposed method's printed value is
better in the metric's direction, and equal printed values count as a tie. This
gives 1,511 comparisons on 520 datasets (one to six metrics per paper, median
three); 16 cannot be decided because a value is not printed as a number. Of the
other 1,495, the proposed method wins 1,107, ties 102 and loses 286, a win in
79.5\% of the decided comparisons (95\% CI 77.3 to 81.5). Twenty-three of the 29
papers report at least one comparison that SMOTE wins.

\begin{table}[t]
\centering\small
\caption{Comparisons of the Proposed Method with Plain SMOTE by Metric (29 Papers).
A row is one dataset and one metric in one paper; printed values are compared in the
direction of the metric, equal printed values are ties, and n.c.\ counts values not
printed as numbers. Scores are not comparable across papers.}
\label{tab:wins-by-metric}
\begin{tabular}{lrrrrrr}
\toprule
Metric & Papers & Rows & Proposed wins & Ties & SMOTE wins & n.c.\\
\midrule
AUC & 19 & 403 & 309 & 30 & 64 & 0\\
G-mean & 19 & 380 & 283 & 29 & 64 & 4\\
F-measure & 18 & 287 & 217 & 16 & 50 & 4\\
Recall (sensitivity) & 12 & 193 & 130 & 21 & 42 & 0\\
Precision & 7 & 113 & 62 & 5 & 46 & 0\\
Accuracy & 4 & 69 & 62 & 1 & 6 & 0\\
MCC & 3 & 52 & 40 & 0 & 8 & 4\\
AUC-PR & 1 & 7 & 3 & 0 & 0 & 4\\
Specificity & 1 & 7 & 1 & 0 & 6 & 0\\
\addlinespace
All metrics & 29 & 1,511 & 1,107 & 102 & 286 & 16\\
\bottomrule
\end{tabular}
\end{table}

Table~\ref{tab:wins-by-metric} shows that the margin depends on the metric. On the
measures that reward finding minority examples (AUC, G-mean, F-measure, recall) the
proposed method wins three to four comparisons in five, while on precision the split
is 62 to 46 and on specificity SMOTE wins six of seven. Counting one result per dataset
(the majority over its metrics) gives 409 wins, 35 ties and 76 losses on 520 datasets,
and the earlier single-metric rule (G-mean, else AUC, else F-measure, else accuracy)
gives 521 comparisons, in which the proposed method wins in 402 rows, ties in 39 and
loses in 80.

\begin{table}[t]
\centering\small
\caption{Comparisons with Plain SMOTE by How the Dataset Was Obtained (All Metrics,
29 Papers). Win share is the share of wins among decided comparisons (ties left out),
with a 95\% Wilson interval. Sixteen comparisons not printed as numbers are left out.}
\label{tab:wins-by-type}
\begin{tabular}{lrrrrrrl}
\toprule
How obtained & Datasets & Papers & Rows & Wins & Ties & Losses & Win share (\%)\\
\midrule
Original & 144 & 25 & 425 & 329 & 10 & 86 & 79.3 (75.1--82.9)\\
Repository-defined & 119 & 9 & 315 & 220 & 44 & 51 & 81.2 (76.1--85.4)\\
Modified & 143 & 14 & 495 & 356 & 20 & 119 & 74.9 (70.9--78.6)\\
Synthetic & 5 & 1 & 15 & 15 & 0 & 0 & 100 (79.6--100)\\
Cannot be told & 105 & 13 & 234 & 184 & 24 & 26 & 87.6 (82.5--91.4)\\
Not matched to a dataset & 4 & 3 & 11 & 3 & 4 & 4 & 42.9 (15.8--75.0)\\
\addlinespace
All & 520 & 29 & 1,495 & 1,107 & 102 & 286 & 79.5 (77.3--81.5)\\
\bottomrule
\end{tabular}
\end{table}

We tested the hypothesis that motivated this analysis: that proposed methods win
more often on data the authors modified or generated than on original data. The data
do not support it (Table~\ref{tab:wins-by-type}). Modified and synthetic data give a
win in 75.7\% of decided comparisons (371 of 490) against 79.3\% on original data.
Because the comparisons of one paper are not independent, we estimated the
difference with the paper as the unit of clustering: a logistic model with
paper-level clustering (generalised estimating equations) gives an odds ratio of 0.68
(95\% CI 0.43 to 1.08, $p = 0.10$) for modified or synthetic against original data;
a conditional logistic model that compares the two kinds of data only inside the 11
papers that use both and report both outcomes gives 0.61 (0.38 to 0.97, $p = 0.04$);
and a mixed model with random effects for paper and for dataset gives 0.75 (0.58 to
0.97). The direction is therefore the opposite of the hypothesis, but it is not
stable: with one result per dataset the two clustered estimates are 0.64 ($p = 0.15$)
and 0.46 ($p = 0.13$), and under the single-metric rule 0.80 ($p = 0.47$) and 0.61
($p = 0.31$). Paper by paper, among the 13 papers with both kinds of data, the win
share is higher on modified or synthetic data in 3, higher on standard data in 7 and
equal in 3 (sign test $p = 0.34$). Repository-defined splits do not differ from
original data (odds ratio 1.19, $p = 0.61$), and the comparisons whose data origin
cannot be told show the highest raw win share, a difference that disappears once
papers are taken into account (odds ratio 0.99, $p = 0.98$).

The second hypothesis concerned papers that add a single standard dataset to data of
another kind. Two papers fit this pattern: MWMOTE (one standard and 19 other
datasets) and RUSBoost~\pcite{seiffert2010rusboost} (one and 14). Both win on their
standard dataset; on their other datasets they win 52 of 65 decided comparisons
(80.0\%), the same share as in the 21 papers with at least two standard datasets
(482 of 600, 80.3\%). Five papers compare with SMOTE only on standard data and one
(INOS) only on modified data. Two papers are too few for a test, and nothing in them
suggests that the standard dataset serves a different purpose from the rest.

Measured difficulty says more about where the wins fall than the origin of the data.
For the 72 verifiable two-class files (Section~\ref{sec:data-difficulty}) that appear
in these comparisons (20 papers, 374 comparisons), the chance that the proposed method
beats SMOTE rises with the noise index: the odds ratio per 0.1 of noise index is 1.15
(95\% CI 1.04 to 1.26, $p = 0.005$) with paper-level clustering and 1.19 (1.04 to 1.35,
$p = 0.01$) within papers, and 1.16 (0.98 to 1.36, $p = 0.08$) with one result per file
and paper. The imbalance ratio shows no such relation (odds ratio 0.92 per doubling,
$p = 0.61$), nor do the small-sample, outlier and small-disjunct flags taken one at a
time. This analysis is exploratory: it covers the files that could be verified, which
are mostly small and low-dimensional.

Comparisons with SMOTE inside another scheme are kept apart and are not pooled with
Tables~\ref{tab:wins-by-metric} and~\ref{tab:wins-by-type}. Twelve papers contribute
558 comparisons: Static-SMOTE (EMDO~\pcite{yao2021emdo}, OVO-SMOTE, SOUP) 79, with 69
wins, 3 ties and 7 losses; SMOTE after one-versus-one decomposition (AMDO~\pcite{yang2018amdo}, GDHS~\pcite{yan2025gdhs},
SHSampler~\pcite{zheng2022shsampler}) 148, with 101, 6 and 41; SMOTE applied to all
classes (MC-CCR~\pcite{koziarski2020mcccr}, MC-RBO~\pcite{krawczyk2020mcrbo}, NROMM) 239,
with 231, 0 and 8; and SMOTE combined with another learner (DSRBF~\pcite{fernandeznavarro2011dsrbf}, EasyEnsemble, GAMO) 92,
with 71, 2 and 19. Three papers are left out of every comparison set. KNNOR~\pcite{islam2022knnor}
prints only means over its 19 datasets (G-score .8363 for KNNOR and .8238 for SMOTE).
MC-MBRC~\pcite{ma2024mcmbrc} names SMOTE as a baseline, but the column headers of its
per-dataset tables do not say which column is SMOTE. In RBBoost~\pcite{diezpastor2015randombalance},
the scores of its RB-Boost come from the Boosting-family table of its supplement
(Supplementary Table~11) and those of its SMOTE ensemble (ESM) from the
data-processing-family table (Supplementary Table~3); both report the same protocol
(J48 with Laplace smoothing, ensembles of 100, stratified $5\times2$-fold
cross-validation), but the paper compares methods only within each family and never
compares RB-Boost with ESM directly. Its 86 datasets, 63 of them repository-defined,
are kept in our data for audit.

\paragraph{What the reader should do.} Read a reported win over SMOTE together with
the metric: most of the margin comes from minority-sensitive measures, and precision
and specificity often favour SMOTE. Expect wins over SMOTE to be more frequent on data
whose minority examples are borderline or isolated, and less informative on clean
data. Whether the data were modified does not, on this evidence, inflate the win
rate; whether they are described well enough to be identified matters more for the
reader (Sections~\ref{sec:data-names} and~\ref{sec:data-verifiability}).

\evidence{Rows: out/winloss\_allmetric\_rows.csv (2,497 dataset-by-metric rows from
ai-draft/outcome/allmetric/*.json, 44 papers; the single-metric values re-copied
independently and identical on 839 of 839 rows, out/winloss\_allmetric\_crosscheck.csv);
build/winloss\_allmetric.py $\rightarrow$ out/winloss\_allmetric\_summary.json (plain,
non-plain groups, exclusions R4, R6, R8; paper classes aligned with
build/reconcile\_counts.py on 29 Oct 2026). Data types: build/allmetric\_links.py
$\rightarrow$ out/allmetric\_links.csv (2,496 of 2,497 rows linked to a record). Tests:
build/hypotheses.py $\rightarrow$ out/hypotheses.json, out/hypotheses\_paper\_split.csv
(statsmodels 0.15: GEE exchangeable, ConditionalLogit, BinomialBayesMixedGLM variational
fit). Single-metric set: out/verified\_outcomes.csv, out/reconcile\_counts.json (521 rows,
rules R4, R6, R8). Difficulty: build/difficulty\_outcomes.py $\rightarrow$
out/difficulty\_outcomes.json. RBBoost: supplement Table~3 pp.3--5 and Table~11
pp.17--18; protocol p.101 and family design p.102. MC-MBRC: SMOTE named in Appendix A.1
and Table~8 only; Tables 9--12, p.10--11.}

